\begin{figure}[htbp]
\centering
\begin{tikzpicture}[>=stealth,scale=1.25]
\draw[->,black!40] (-1.4,0)--(4.2,0);
\draw[->,black!40] (0,-1)--(0,2.7);
\draw[->,arkivblaa] (0,0)--(1.4,.8) node[above left] {$v$};
\draw[->,black!65,dashed] (0,0)--(2.8,1.6) node[above] {$2v$};
\draw[->,black!65,dotted] (0,0)--(-.7,-.4) node[below] {$-\frac12v$};
\end{tikzpicture}
\caption{Skalarmultiplikation ændrer længden med $|a|$ og vender retningen ved $a<0$; $a=0$ giver nulvektoren. Her har de to tegnede retninger samme dimensionsløse målestok. Vektoraddition vises ved parallelogrammet i basisfiguren.}
\label{fig:b1-vektorskalering}
\end{figure}
